\chapter{Brain}

\section*{Patient focus}
This chapter explains the examination or clinical problem in accessible language. Patients should use it to prepare questions and understand the purpose, experience, benefits, and limitations of imaging; local instructions and the official report take priority.


\label{ch:ch38_brain}

When your doctor orders a brain scan, the most common options are CT (computed tomography) and MRI (magnetic resonance imaging). A CT scan uses X-rays and takes only a few minutes. It is very good at detecting bleeding in the brain and is almost always the first test done in an emergency, such as a sudden severe headache or after a head injury. An MRI uses strong magnets and radio waves. It takes longer (30--60 minutes) but produces much more detailed pictures of the brain tissue, nerves, and blood vessels. Neither test is painful. For CT, you may receive an injection of contrast dye through a vein; for MRI, the machine is noisy and you must lie very still inside a narrow tube.

Brain imaging can detect many conditions: bleeding (which may be outside or inside the brain tissue), stroke, tumours, infections such as abscesses, swelling, and problems with the fluid-filled spaces called ventricles. The results help your doctor decide whether you need medication, surgery, or further tests.

---

\section{Imaging Strategy}
Select CT or MRI according to urgency, suspected pathology, stability, and availability. CT is fast for hemorrhage, fracture, and acute mass effect; MRI better characterizes parenchyma, posterior fossa, cranial nerves, and many tumors. CTA, CTV, perfusion, or contrast may be added for a defined question.

\begin{table}[h]
\centering
\caption{CT vs MRI for brain imaging}
\begin{tabular}{p{3.5cm} p{5.5cm} p{5.5cm}}
\hline
\textbf{Feature} & \textbf{CT} & \textbf{MRI} \\
\hline
Acute hemorrhage & Excellent (hyperdense) & Excellent (variable signal) \\
Bone detail & Excellent & Limited \\
Contraindications & Radiation, contrast allergy & Claustrophobia, metallic implants, pacemakers \\
Cranial nerves & Poor & Excellent \\
Perfusion & CTP available & PWI available \\
Posterior fossa & Beam-hardening artifact & Superior \\
Speed & 2--5 minutes & 30--60 minutes \\
White matter disease & Limited & Superior (FLAIR) \\
\hline
\end{tabular}
\end{table}

\begin{pearlbox}[CT first for acute presentations]
In any acute neurological emergency---trauma, sudden-onset severe headache, or focal deficit---non-contrast CT is the first-line investigation. It excludes hemorrhage in under 5 minutes and determines whether thrombolysis is safe.
\end{pearlbox}

---

\section{Hemorrhage Patterns}
Intracranial hemorrhage is classified by location, each with characteristic imaging features and clinical associations.

\begin{table}[h]
\centering
\caption{Types of intracranial hemorrhage}
\begin{tabular}{p{2.8cm} p{3.5cm} p{3.5cm} p{3.5cm}}
\hline
\textbf{Type} & \textbf{Location} & \textbf{CT appearance} & \textbf{Common causes} \\
\hline
Epidural & Between skull and dura & Biconvex (lens-shaped), does not cross sutures & Middle meningeal artery trauma \\
Intraparenchymal & Within brain tissue & Round/oval hyperdensity & Hypertension, amyloid angiopathy, trauma \\
Intraventricular & Within ventricles & Fluid--blood levels, fills ventricles & Extension from other hemorrhage, AV malformation \\
Subarachnoid & Within subarachnoid space & Blood in sulci, cisterns, fissures & Ruptured aneurysm (85\%), trauma \\
Subdural & Between dura and arachnoid & Crescentic, crosses sutures, may cross midline & Bridging vein tears (elderly, anticoagulation) \\
\hline
\end{tabular}
\end{table}

\begin{pearlbox}[Epidural vs subdural on CT]
Epidural hematomas are biconvex and respect suture lines (the dura is firmly attached at sutures). Subdural hematomas are crescentic and can spread along the entire hemisphere because the dura is not attached at sutures. This distinction is critical because epidural hematomas often require emergency surgical evacuation.
\end{pearlbox}

---

\section{Stroke Imaging}
Acute stroke imaging follows a structured protocol:

\begin{enumerate}
  \item \textbf{Non-contrast CT (NCCT)}: Excludes hemorrhage (mandatory before thrombolysis). Look for early signs: loss of insular ribbon, hyperdense MCA sign, effacement of sulci.
  \item \textbf{CT Angiography (CTA)}: Identifies large vessel occlusion (LVO) for endovascular thrombectomy. Examine the internal carotid, M1/M2 MCA, and basilar arteries.
  \item \textbf{CT Perfusion (CTP)}: Quantifies infarct core (CBF $< 30\%$) and penumbra (Tmax $> 6$ s). A large mismatch extends the treatment window to 24 hours.
  \item \textbf{MRI with DWI}: Diffusion-weighted imaging detects cytotoxic edema within minutes, appearing as restricted diffusion (bright on DWI, dark on ADC). DWI--FLAIR mismatch confirms hyperacute onset ($< 4.5$ h).
\end{enumerate}

See Chapter~\ref{ch:stroke_imaging} for detailed stroke imaging protocol.

---

\section{Tumor Imaging}
Brain tumors are characterized by location, enhancement pattern, edema, and mass effect.

\begin{table}[h]
\centering
\caption{Common brain tumors and imaging features}
\begin{tabular}{p{2.5cm} p{3cm} p{3cm} p{3cm}}
\hline
\textbf{Tumor} & \textbf{Location} & \textbf{CT} & \textbf{MRI} \\
\hline
Glioblastoma & Cerebral hemispheres, crosses midline & Heterogeneous mass, ring enhancement, central necrosis & T2/FLAIR hyperintense, ring enhancement, restricted diffusion at rim \\
Low-grade glioma & Frontal/temporal lobes & Hypodense, no enhancement & T2/FLAIR hyperintense, no enhancement \\
Meningioma & Dural-based, parasagittal, sphenoid wing & Hyperdense, homogeneous enhancement, calcification, hyperostosis & T1 isointense, intense homogeneous enhancement, dural tail \\
Metastases & Grey--white junction, multiple & Ring-enhancing, surrounding edema & T1 ring enhancement, extensive vasogenic edema \\
Pituitary adenoma & Sellar region & Enlarged sella, enhancement & Macroadenoma $> 10$ mm, chiasmal compression \\
\hline
\end{tabular}
\end{table}

\begin{pearlbox}[Multiple ring-enhancing lesions]
The differential for multiple ring-enhancing brain lesions includes metastases (most common in adults), abscesses (restricted diffusion centrally), high-grade glioma, lymphoma, and demyelination (open ring sign). Clinical context and DWI pattern are key to narrowing the differential.
\end{pearlbox}

---

\section{Infection}

\subsection{Brain Abscess}
A brain abscess appears as a ring-enhancing lesion with surrounding vasogenic edema and mass effect. The pus within the abscess shows \textbf{restricted diffusion} (bright on DWI, dark on ADC), which helps distinguish it from necrotic tumor (which typically does not restrict centrally). Common sources include sinusitis, otitis media, dental infection, endocarditis, or hematogenous spread.

\subsection{Encephalitis}
Herpes simplex encephalitis characteristically involves the medial temporal lobes and insular cortices with T2/FLAIR hyperintensity and restricted diffusion. It is a neurological emergency requiring immediate IV aciclovir. Other causes include autoimmune (anti-NMDA receptor) and other viral encephalitides.

---

\section{Hydrocephalus}
Hydrocephalus refers to enlargement of the cerebral ventricles due to impaired CSF flow or absorption.

\begin{table}[h]
\centering
\caption{Types of hydrocephalus}
\begin{tabular}{p{3cm} p{5cm} p{5cm}}
\hline
\textbf{Type} & \textbf{Mechanism} & \textbf{Key features} \\
\hline
Communicating & Impaired CSF absorption at arachnoid granulations & All ventricles enlarged uniformly \\
Normal pressure & Chronic, often in elderly; impaired absorption & Ventriculomegaly out of proportion to sulcal atrophy; episodic gait disturbance, dementia, incontinence \\
Obstructive (non-communicating) & Blockage within ventricular system & Dilated ventricles proximal to obstruction; e.g.\ aqueductal stenosis \\
\hline
\end{tabular}
\end{table}

\begin{pearlbox}[Transependymal edema]
On CT or MRI, periventricular hypodensity (CT) or T2 hyperintensity (MRI) surrounding the ventricles indicates transependymal edema---CSF being forced through the ependyma into the brain parenchyma. This is a sign of acute or subacute hydrocephalus requiring urgent evaluation.
\end{pearlbox}

---

\section{Structured Report: Brain Hemorrhage}

\begin{testbox}[Indications:]
68-year-old male on warfarin presenting with sudden left-sided weakness and decreased consciousness. INR 3.8.
\end{testbox}
\begin{refbox}[Findings:]
\textbf{Non-contrast CT head}: There is an acute right-sided subdural hematoma measuring up to 12 mm in maximum thickness, extending over the right frontoparietal convexity. The hematoma is crescentic and crosses suture lines. There is midline shift of 8 mm to the left at the level of the septum pellucidum. Effacement of the right lateral ventricle is noted. No intraparenchymal hemorrhage or subarachnoid hemorrhage. The basal cisterns are partially effaced. No skull fracture.
\end{refbox}
\begin{pearlbox}[Impression:]
1. Acute right frontoparietal subdural hematoma, 12 mm maximum thickness, with 8 mm midline shift---\textbf{urgent neurosurgical review indicated}.
2. Mass effect with effacement of the right lateral ventricle and partial basal cistern effacement.
3. Consider urgent hematoma evacuation and reversal of anticoagulation.
\end{pearlbox}

---

\section{Systematic Review}
Assess scalp and skull, extra-axial spaces, ventricles, cisterns, gray-white differentiation, deep nuclei, white matter, posterior fossa, sellar region, vascular structures, and diffusion or susceptibility when available. Compare prior imaging and state technical limitations.

\section{Safety}
Urgent hemorrhage, herniation, hydrocephalus, large-vessel occlusion, and acute mass effect require prompt communication. MRI screening and contrast decisions are patient- and device-specific.

\clearpage
